\documentclass[11pt]{article}
\usepackage[margin=0.85in]{geometry}
\usepackage{amsmath,amssymb,microtype}
\usepackage[T1]{fontenc}
\usepackage{lmodern}
\usepackage[hidelinks]{hyperref}
\title{Two Affirmative Answers for Continuous Operator Families}
\author{Literature-implied status note for arXiv:2010.07040}
\date{August 11, 2026}

\begin{document}
\maketitle

\section*{Status and source questions}
\textbf{Status: literature-implied answer (both questions, affirmative).}

After Theorem~1.5, Berkani asks the following two questions
\cite[p.~4]{Berkani}.
\begin{enumerate}
\item Given $S\in C(\mathbb X,L(X)/K(X))$, must there be
$T\in C(\mathbb X,L(X))$ whose pointwise Calkin image is $S$?
\item If $H$ is a Hilbert space and
$\mathcal K\in C(\mathbb X,K(H))$, is $\mathcal K$ a uniform norm limit of
continuous finite-rank families?
\end{enumerate}
Here $\mathbb X$ is compact and all function spaces carry the supremum norm.

\section*{Question 1: continuous lifting}
Let
\[
 q:L(X)\longrightarrow L(X)/K(X)
\]
be the quotient map.  It is a continuous linear surjection between Banach
spaces.  The classical Bartle--Graves theorem says that every such map has a
continuous, positively homogeneous right inverse
\cite[Theorem~5.1, p.~14]{Messerschmidt}; thus there is a continuous
$s:L(X)/K(X)\to L(X)$ with $q\circ s=\mathrm{id}$.

For the given family $S$, define
\[
 T_x=s(S_x)\qquad(x\in\mathbb X).
\]
Then $T=s\circ S$ is continuous and $q(T_x)=S_x$ for every $x$.  Since
$\mathbb X$ is compact, $T$ is automatically bounded and belongs to the
source's Banach algebra $C(\mathbb X,L(X))$.  Therefore the natural map
\[
 C(\mathbb X,L(X))/C(\mathbb X,K(X))
 \longrightarrow C(\mathbb X,L(X)/K(X))
\]
is surjective.  In fact, the lifting conclusion itself does not require
compactness of the parameter space.

\section*{Question 2: one compression for the whole family}
Let
\[
 \mathcal C_0=\{\mathcal K_x:x\in\mathbb X\}\subset K(H).
\]
This is compact because it is the continuous image of the compact space
$\mathbb X$.  Fix $\varepsilon>0$.  Choose a finite
$\varepsilon/4$-net $K^{(1)},\ldots,K^{(m)}$ in $\mathcal C_0$.  Since
finite-rank operators are norm dense in $K(H)$, choose finite-rank $F_j$ with
\[
 \|K^{(j)}-F_j\|<\varepsilon/4.
\]
Let $E$ be the finite-dimensional span of all
$\operatorname{ran}F_j$ and $\operatorname{ran}F_j^*$, and let $P$ be the
orthogonal projection onto $E$.  The two range inclusions give
$PF_j=F_j$ and $F_jP=F_j$, hence $PF_jP=F_j$.

For any $x$, select $j$ with
$\|\mathcal K_x-K^{(j)}\|<\varepsilon/4$.  Then
\[
 \|\mathcal K_x-F_j\|<\varepsilon/2
\]
and, because $P$ is contractive,
\[
 \|\mathcal K_x-P\mathcal K_xP\|
 \leq \|\mathcal K_x-F_j\|+
       \|P(F_j-\mathcal K_x)P\|
 <\varepsilon.
\]
The family $x\mapsto P\mathcal K_xP$ is norm-continuous and every member has
range in $E$, so it is a continuous finite-rank family.  Taking
$\varepsilon=1/n$ produces the requested uniformly convergent sequence.  No
separability assumption on $H$ is needed.

\section*{Provenance and scope}
The first answer is a direct specialization of the Bartle--Graves theorem,
which long predates the source question.  The second is the standard density
of finite-rank operators in $K(H)$ made uniform over a compact subset by a
single finite-dimensional compression.  The relation to Berkani's two
questions is agent-identified: the supporting authors do not claim to answer
them.

A bounded search through August 11, 2026 used the exact arXiv id and title,
the question wording, continuous Calkin lift terms, Bartle--Graves, and
continuous finite-rank-family terms.  It found the source, the supporting
selection theorem, and general uses of continuous Calkin sections, but no
later paper explicitly presenting itself as a resolution of these two
questions.

\begin{thebibliography}{9}
\bibitem{Berkani}
M. Berkani,
\emph{Index of continuous families of bounded linear operators in Banach spaces and application},
arXiv:2010.07040 (2020),
\url{https://arxiv.org/abs/2010.07040}.

\bibitem{Messerschmidt}
M. Messerschmidt,
\emph{A pointwise Lipschitz selection theorem},
arXiv:1611.08435; see Theorem~5.1 for the classical Bartle--Graves theorem,
\url{https://arxiv.org/abs/1611.08435}.

\bibitem{BartleGraves}
R.~G. Bartle and L.~M. Graves,
\emph{Mappings between function spaces},
Trans. Amer. Math. Soc. \textbf{72} (1952), 400--413.
\end{thebibliography}

\end{document}
